\section{Introduction}

Robotic manipulation involves intricate interactions with objects in the environment, which can often be expressed as constraints in both spatial and temporal domains.
Consider the task of pouring tea into a cup in Fig.~\ref{fig:teaser}: the robot must grasp \emph{at the handle}, keep the cup \emph{upright} while transporting it, \emph{align} the spout with the target container, and then tilt the cup \emph{at the correct angle} to pour.
Here, the constraints encode not only the intermediate sub-goals (e.g., align the spout) but also the transitioning behaviors (e.g., keep the cup \emph{upright} in transportation), which collectively dictate the spatial, timing, and other combinatorial requirements of the robot's actions in relation to the environment.

However, effectively formulating these constraints for a large variety of real-world tasks presents significant challenges.
While representing constraints using relative poses between robots and objects is a direct and widely-used approach~\cite{kaelbling2010hierarchical}, rigid-body transformations do not depict geometric details, require obtaining object models \textit{a priori}, and cannot work on deformable objects.
On the other hand, data-driven approaches enable learning constraints directly in visual space~\cite{driess2022learning,simeonov2022neural}. While more flexible, it remains unclear how to effectively collect training data as the number of constraints grows combinatorially in terms of objects and tasks.
Therefore, we ask the question: how can we represent constraints in manipulation that are 1) widely applicable: adaptable to tasks that require multi-stage, in-the-wild, bimanual, and reactive behaviors, 2) scalably obtainable: have the potential to be fully automated through the advances in foundation models, and 3) real-time optimizable: can be efficiently solved by off-the-shelf solvers to produce complex manipulation behaviors?

In this work, we propose~\textbf{\algname~(\algabrvname)}.
Specifically,~\algabrvname represents constraints as Python functions that map a set of keypoints to a numerical cost, where each keypoint is a task-specific and semantically meaningful 3D point in the scene.
Each function is composed of (potentially nonlinear) arithmetic operations on the keypoints and encodes a desired ``relation'' between them, where the keypoints may belong to different entities in the environment, such as the robot arms, object parts, and other agents.
While each keypoint only consists of its 3D Cartesian coordinates in the world frame, multiple keypoints can collectively specify lines, surfaces, and/or 3D rotations if rigidity between keypoints is enforced.
We study~\algabrvname in the context of the sequential manipulation problem, where each task involves multiple stages that have spatio-temporal dependencies (e.g., ``grasping'', ``aligning'', and ``pouring'' in the aforementioned example).



While constraints are typically defined manually per task~\cite{manuelli2019kpam},
we demonstrate the specific form of~\algabrvname possesses a unique advantage in that they can be automated by pre-trained large vision models (LVM)~\cite{oquab2023dinov2} and vision-language models (VLM)~\cite{openai2023gpt}, enabling \emph{in-the-wild} specification of~\algabrvname from RGB-D observations and free-form language instructions.
Specifically, we leverage LVM to propose fine-grained and semantically meaningful keypoints in the scene and VLM to write the constraints as Python functions from visual input overlaid with proposed keypoints.
This process can be interpreted as grounding fine-grained spatial relations, often those not easily specified with natural language, in an output modality supported by VLM (code) using visual referral expressions.

With the generated constraints, off-the-shelf solvers can be used to produce robot actions by re-evaluating the constraints based on tracked keypoints. Inspired by~\cite{toussaint2022sequence}, we employ a hierarchical optimization procedure to first solve a set of waypoints as sub-goals (represented as $SE(3)$ end-effector poses) and then solve the receding-horizon control problem to obtain a dense sequence of actions to achieve each sub-goal.
With appropriate instantiation of the problem, we demonstrate that it can be reliably solved at approximately 10 Hz for the tasks considered in this work.


Our contributions are summarized as follows: 1) We formulate manipulation tasks as a hierarchical optimization problem with~\algname; 2) We devise a pipeline to automatically specify keypoints and constraints using large vision models and vision-language models; 3) We present system implementations on two real-robot platforms that take as input a language instruction and RGB-D observations, and produce multi-stage, in-the-wild, bimanual, and reactive behaviors for a large variety of manipulation tasks, all without task-specific data or environment models. 
